{\large\textbf{Hydrogen and galaxies (10 points)}}

This problem aims to study the peculiar physics of galaxies, such as their
dynamics and structure. In particular, we explain how to measure the mass
distribution of our galaxy from the inside. For this we will focus on hydrogen,
its main constituent.

Throughout this problem we will only use $\hbar$, defined as $\hbar = h/2\pi$.

\textbf{Part A - Introduction}

\textbf{Bohr model}

We assume that the hydrogen atom consists of a non-relativistic electron, with
mass $m_e$, orbiting a fixed proton. Throughout this part, we assume its motion
is on a circular orbit.

\textbf{A.1}\quad Determine the electron's velocity $v$ in a circular orbit of
radius $r$. \hfill 0.2pt

In the Bohr model, we assume the magnitude of the electron's angular momentum
$L$ is quantized, $L = n\hbar$ where $n > 0$ is an integer. We define
$\alpha = \frac{e^2}{4\pi\varepsilon_0\hbar c} \approx 7.27 \times 10^{-3}$.

\textbf{A.2}\quad Show that the radius of each orbit is given by
$r_n = n^2 r_1$, where $r_1$ is called the Bohr radius. Express $r_1$ in terms
of $\alpha$, $m_e$, $c$ and $\hbar$ and calculate its numerical value with 3
digits. Express $v_1$, the velocity on the orbit of radius $r_1$, in terms of
$\alpha$ and $c$. \hfill 0.5pt

\textbf{A.3}\quad Determine the electron's mechanical energy $E_n$ on an orbit
of radius $r_n$ in terms of $e$, $\varepsilon_0$, $r_1$ and $n$. Determine $E_1$
in the ground state in terms of $\alpha$, $m_e$ and $c$. Compute its numerical
value in eV. \hfill 0.5pt

\textbf{Hydrogen fine and hyperfine structures}

The rare spontaneous inversion of the electron's spin causes a photon to be
emitted on average once per 10 million years per hydrogen atom. This emission
serves as a hydrogen tracer in the universe and is thus fundamental in
astrophysics. We will study the transition responsible for this emission in two
steps.

First, consider the interaction between the electron spin and the relative
motion of the electron and the proton. Working in the electron's frame of
reference, the proton orbits the electron at a distance $r_1$. This produces a
magnetic field $\overrightarrow{B}_1$.

\textbf{A.4}\quad Determine the magnitude $B_1$ of $\overrightarrow{B}_1$ at the
position of the electron in terms of $\mu_0$, $e$, $\alpha$, $c$ and $r_1$.
\hfill 0.5pt

% script M printed as ℳ; \mathcal{M} is the closest symbol in the allowed packages
Second, the electron spin creates a magnetic moment
$\overrightarrow{\mathcal{M}}_s$. Its magnitude is roughly
$\mathcal{M}_s = \frac{e}{m_e}\hbar$. The \textit{fine} (F) structure is related
to the energy difference $\varDelta E_\mathrm{F}$ between an electron with a
magnetic moment $\overrightarrow{\mathcal{M}}_s$ parallel to
$\overrightarrow{B}_1$ and that of an electron with
$\overrightarrow{\mathcal{M}}_s$ anti-parallel to $\overrightarrow{B}_1$.
Similarly, the \textit{hyperfine} (HF) structure is related to the energy
difference $\varDelta E_\mathrm{HF}$, due to the interaction between parallel
and anti-parallel magnetic moments of the electron and the proton. It is known
to be approximately
$\varDelta E_\mathrm{HF} \simeq 3.72\frac{m_e}{m_p}\varDelta E_\mathrm{F}$
where $m_p$ is the proton mass.

\textbf{A.5}\quad Express $\varDelta E_\mathrm{F}$ as a function of $\alpha$ and
$E_1$.\\
Express the wavelength $\lambda_\mathrm{HF}$ of a photon emitted during a
transition between the two states of the hyperfine structure and give its
numerical value with two digits. \hfill 0.5pt

\textbf{Part B - Rotation curves of galaxies}

\textbf{Data}
\begin{itemize}
\item Kiloparsec: $1\,\mathrm{kpc} = 3.09 \times 10^{19}\,\mathrm{m}$
\item Solar mass : $1\,\mathrm{M}_\odot = 1.99 \times 10^{30}\,\mathrm{kg}$
\end{itemize}

We consider a spherical galaxy centered around a fixed point $O$. At any point
$P$, let $\rho = \rho(P)$ be the volumetric mass density and
$\varphi = \varphi(P)$ the associated gravitational potential (i.e.\ potential
energy per unit mass). Both $\rho$ and $\varphi$ depend only on
$r = \left\|\overrightarrow{OP}\right\|$. The motion of a mass $m$ located at
$P$, due to the field $\varphi$, is restricted to a plane containing $O$.

\textbf{B.1}\quad In the case of a circular orbit, determine the velocity $v_c$
of an object on a circular orbit passing through $P$ in terms of $r$ and
$\dfrac{d\varphi}{dr}$. \hfill 0.2pt

Fig.~1(A) is a picture of the spiral galaxy NGC 6946 in the visible band (from
the $0.8\,\mathrm{m}$ Schulman Telescope at the Mount Lemmon Sky Center in
Arizona). The little ellipses in Fig.~1(B) show experimental measurements of
$v_c$ for this galaxy. The central region ($r < 1\,\mathrm{kpc}$) is named the
bulge. In this region, the mass distribution is roughly homogeneous. The red
curve is a prediction for $v_c$ if the system were homogeneous in the bulge and
keplerian ($\varphi(r) = -\beta/r$ with $\beta > 0$) outside it, i.e.\
considering that the total mass of the galaxy is concentrated in the bulge.

\begin{center}\includegraphics[width=0.86\linewidth]{figures/IPhO_2025_Q1_fig1.png}\end{center}

\begin{center}\textbf{Fig.~1}: NGC 6946 galaxy: Picture (A) and rotation curve (B).\end{center}

\textbf{B.2}\quad Deduce the mass $M_b$ of the bulge of NGC 6946 from the red
rotation curve in Fig.~1(B), in solar mass units. \hfill 0.5pt

Comparing the keplerian model and the experimental data makes astronomers
confident that part of the mass is invisible in the picture. They thus suppose
that the galaxy's actual mass density is given by
\begin{equation}
\rho_m(r) = \frac{C_m}{r_m^2 + r^2} \tag{1}
\end{equation}
where $C_m > 0$ and $r_m > 0$ are constants.

\textbf{B.3}\quad Show that the velocity profile $v_{c,m}(r)$, corresponding to
the mass density in Eq.~1, can be written
$v_{c,m}(r) = \sqrt{k_1 - \frac{k_2\cdot\arctan(\frac{r}{r_m})}{r}}$. Express
$k_1$ and $k_2$ in terms of $C_m$, $r_m$ and $G$.\\
( Hints: $\displaystyle\int_0^r \frac{x^2}{a^2 + x^2} dx = r - a\arctan(r/a)$,
and: $\arctan(x) \simeq x - x^3/3$ for $x \ll 1$. )

Simplify $v_{c,m}(r)$ when $r \ll r_m$ and when $r \gg r_m$.\\
Show that if $r \gg r_m$, the mass $M_m(r)$ embedded in a sphere of radius $r$
with the mass density given by Eq.~1 simplifies and depends only on $C_m$ and
$r$.\\
Estimate the mass of the galaxy NGC 6946 actually present in the picture in
Fig.~1(A). \hfill 1.8pt

\textbf{Part C - Mass distribution in our galaxy}

For a spiral galaxy, the model for Eq.~1 is modified and one usually considers
the gravitational potential is given by
$\varphi_G(r,z) = \varphi_0 \ln\left(\dfrac{r}{r_0}\right)
\exp\left[-\left(\dfrac{z}{z_0}\right)^2\right]$, where $z$ is the distance to
the galactic plane (defined by $z = 0$ ), and $r < r_0$ is now the axial radius
and $\varphi_0 > 0$ a constant to be determined. $r_0$ and $z_0$ are constant
values.

\textbf{C.1}\quad Find the equation of motion on $z$ for the vertical motion of
a point mass $m$ in such a potential, assuming $r$ is constant. Show that, if
$r < r_0$, the galactic plane is a stable equilibrium state by giving the
angular frequency $\omega_0$ of small oscillations around it. \hfill 0.5pt

From here on, we set $z = 0$.

\textbf{C.2}\quad Identify the regime, either $r \gg r_m$ or $r \ll r_m$, in
which the model of Eq.~1 recovers a potential of the form $\varphi_G(r,0)$ with
a suitable definition of $\varphi_0$.\\
Under this condition $v_c(r)$ no longer depends on $r$. Express it in terms of
$\varphi_0$. \hfill 0.6pt

Therefore, outside the bulge the velocity modulus $v_c$ does not depend on the
distance to the galactic center. We will use this fact, as astronomers do, to
measure the galaxy's mass distribution from the inside.

All galactic objects considered here for astronomical observations, such as
stars or nebulae, are primarily composed of hydrogen. Outside the bulge, we
assume that they rotate on circular orbits around the galactic center $C$. $S$
is the sun's position and $E$ that of a given galactic object emitting in the
hydrogen spectrum. In the galactic plane, we consider a line of sight $SE$
corresponding to the orientation of an observation, on the unit vector
$\hat{u}_v$ (see Fig.~2).

\begin{center}\includegraphics[width=0.62\linewidth]{figures/IPhO_2025_Q1_fig2.png}\end{center}

\begin{center}\textbf{Fig.~2}: Geometry of the measurement\end{center}

Let $\ell$ be the galactic longitude, measuring the angle between $SC$ and the
$SE$. The sun's velocity on its circular orbit of radius
$R_\odot = 8.00\,\mathrm{kpc}$ is denoted $\overrightarrow{v}_\odot$. A galactic
object in $E$ orbits on another circle of radius $R$ at velocity
$\overrightarrow{v}_E$. Using a Doppler effect on the previously studied
$21\,\mathrm{cm}$ line, one can obtain the relative radial velocity $v_{rE/S}$
of the emitter $E$ with respect to the sun $S$ : it is the projection of
$\vec{v}_E - \vec{v}_\odot$ on the line of sight.

\textbf{C.3}\quad Determine $v_{rE/S}$ in terms of $\ell$, $R$, $R_\odot$ and
$v_\odot$. Then, express $R$ in terms of $R_\odot$, $v_\odot$, $\ell$ and
$v_{rE/S}$. \hfill 0.7pt

Using a radio telescope, we make observations in the plane of our galaxy toward
a longitude $\ell = 30^\circ$. The frequency band used contains the
$21\,\mathrm{cm}$ line, whose frequency is $f_0 = 1.42\,\mathrm{GHz}$. The
results are reported in Fig.~3.

\begin{center}\includegraphics[width=0.61\linewidth]{figures/IPhO_2025_Q1_fig3.png}\end{center}

\textbf{Fig.~3}: Electromagnetic signal as a function of the frequency shift,
measured in the radio frequency band at $\ell = 30^\circ$ using EU-HOU
RadioAstronomy

\textbf{C.4}\quad In our galaxy, $v_\odot = 220\,\mathrm{km\cdot s^{-1}}$.
Determine the values of the relative radial velocity (with 3 significant
digits) and the distance from the galactic center (with 2 significant digits)
of the 3 sources observed in Fig.~3. Distances should be expressed as multiples
of $R_\odot$. \hfill 0.6pt

\textbf{C.5}\quad On the top view of our galaxy (in the answer box), indicate
the positions of the sources observed in Fig.~3.\\
What could be deduced from repeated measurements changing $\ell$? \hfill 0.6pt

\textbf{Part D - Tully-Fisher relation and MOND theory}

The flat external velocity curve of NGC 6946 in Fig.~1 is a common property of
spiral galaxies, as can be seen in Fig.~4 (left). Plotting the external
constant velocity value $v_{c,\infty}$ as a function of the measured total mass
$M_\mathrm{tot}$ of each galaxy gives an interesting correlation called the
Tully-Fischer relation, see Fig.~4 (right).

\begin{center}\includegraphics[width=0.9\linewidth]{figures/IPhO_2025_Q1_fig4.png}\end{center}

\textbf{Fig.~4}. Left: Rotation curves for typical spiral galaxies - Right:
$\log_{10}(M_\mathrm{tot})$ as a function of $\log_{10}(v_{c,\infty})$ on
linear scales. Colored dots correspond to different galaxies and different
surveys. The green line is the Tully-Fischer relation which is in very good
agreement with the best fit line of the data (in black).

\textbf{D.1}\quad Assuming that the radius $R$ of a galaxy doesn't depend on its
mass, show that the model of Eq.~1 (part B) gives a relation of the form
$M_\mathrm{tot} = \eta v_{c,\infty}^\gamma$ where $\gamma$ and $\eta$ should be
specified.\\
Compare this expression to the Tully-Fischer relation by computing
$\gamma_{TF}$. \hfill 0.4pt

In the extremely low acceleration regime, of the order of
$a_0 = 10^{-10}\,\mathrm{m\cdot s^{-2}}$, the MOdified Newtonian Dynamics (MOND)
theory suggests that one can modify Newton's second law using
$\overrightarrow{F} = m\mu\left(\dfrac{a}{a_0}\right)\overrightarrow{a}$ where
$a = \left\|\overrightarrow{a}\right\|$ is the modulus of the acceleration and
the $\mu$ function is defined by $\mu(x) = \dfrac{x}{1 + x}$.

\textbf{D.2}\quad Using data for NGC 6946 in Fig.~1, estimate, within Newton's
theory, the modulus of the acceleration $a_m$ of a mass in the outer regions of
NGC 6946. \hfill 0.2pt

\textbf{D.3}\quad Let $m$ be a mass on a circular orbit of radius $r$ with
velocity $v_{c,\infty}$ in the gravity field of a fixed mass $M$.\\
Within the MOND theory, with $a \ll a_0$, determine the Tully-Fischer
exponent.\\
Using data for NGC 6946 and/or Tully-Fischer law, calculate $a_0$ to show that
MOND operates in the correct regime. \hfill 0.8pt

\textbf{D.4}\quad Considering relevant cases, determine $v_c(r)$ for all values
of $r$ in the MOND theory in the case of a gravitational field due to a
homogeneously distributed mass $M$ with radius $R_b$. \hfill 0.9pt
